\section{Known Design Tensions}
\label{sec:DesignTensions}

Durga sits in a deliberate set of tensions. These are not accidental shortcomings; they are the
result of choosing explicitness and inspectability over hiding complexity behind an engine.

\subsection{Explicitness versus compactness}

The generated runtime is larger and more verbose than a hidden workflow runtime would be. That
cost is accepted because the system is easier to inspect, teach and debug when the BPMN-to-Kafka
mapping is visible.

\subsection{BPMN fidelity versus pragmatic Kafka mapping}

Durga aims to support useful BPMN semantics, but it does so through explicit Kafka-oriented
handlers and topics rather than through a full traditional BPMN engine. That means some semantics
are intentionally pragmatic rather than exhaustive.

\subsection{Regeneration versus local customization}

Generated projects invite local edits, but regeneration remains important. This creates a normal
tension between treating the output as disposable scaffolding and treating it as an evolving code
base. The project currently resolves that by keeping the BPMN model primary and the generated
runtime explicit.

\subsection{Exploration versus production hardening}

The repository is already strong enough for serious architectural exploration, but not every part
of the system is optimized for hardened production operation. The current center of gravity remains
clarity, runnable demos, and generator correctness.
